\section{第12章の解答}\label{解答:第12章}

\paragraph{12.1}
$f(x)=\log(1+x)$の四階微分は
\[
 f^{(4)}(t)=-\frac6{(1+t)^4}
\]
である.\,$0\leqq t\leqq x\leqq\frac12$では$-6\leqq f^{(4)}(t)\leqq0$なので,三次の積分剰余は
\[
 -\int_0^x(x-t)^3\dd t
 \leqq R_3(x)=\frac16\int_0^x(x-t)^3f^{(4)}(t)\dd t
 \leqq0.
\]
左辺は$-\frac{x^4}{4}$である.\,$\log(1+x)=P(x)+R_3(x)$に戻すと
\[
 0\leqq P(x)-\log(1+x)\leqq\frac{x^4}{4}
\]
を得る.\,$x=\frac1{10}$では
\[
 P\left(\frac1{10}\right)
 =\frac1{10}-\frac1{200}+\frac1{3000}
 =\frac{143}{1500},\qquad
 \frac{x^4}{4}=\frac1{40000}.
\]
したがって,求める区間の一つは
\[
 \boxed{\frac{143}{1500}-\frac1{40000}
 \leqq\log\frac{11}{10}\leqq\frac{143}{1500}}
\]
である.両端は有理数であり,区間の幅はちょうど$\frac1{40000}$である.

\paragraph{12.2}
$g(u)=\e^{-u}$について$g'''(u)=-\e^{-u}$なので,\,$u\geqq0$では二次の積分剰余が負になる.さらに$|g'''(u)|\leqq1$より
\[
 \e^{-u}=1-u+\frac{u^2}{2}+r_2(u),
 \qquad -\frac{u^3}{6}\leqq r_2(u)\leqq0.
\]
$u=x^2$を代入して$0$から$\frac12$まで積分すると,近似値は
\[
 Q=\int_0^{\frac12}\left(1-x^2+\frac{x^4}{2}\right)\dd x
 =\frac12-\frac1{24}+\frac1{320}
 =\frac{443}{960}.
\]
真の値を$A$とすれば
\[
 0\leqq Q-A
 \leqq\int_0^{\frac12}\frac{x^6}{6}\dd x
 =\frac1{42}\left(\frac12\right)^7
 =\frac1{5376}.
\]
したがって$\frac{443}{960}$は真の値を上から近似し,誤差は$\frac1{5376}$以下である.

\paragraph{12.3}
分割幅は$h=\frac12$であり,分点$0,\frac12,1,\frac32,2$での関数値は
$0,\frac1{16},1,\frac{81}{16},16$である.これをそれぞれの公式に代入すると
\[
 \begin{aligned}
 T_4&=\frac12\left(8+\frac1{16}+1+\frac{81}{16}\right)
 =\frac{113}{16},\\
 S_4&=\frac16\left\{16+4\left(\frac1{16}+\frac{81}{16}\right)+2\right\}
 =\frac{77}{12}.
 \end{aligned}
\]
真の値は$A=\left[\frac{x^5}{5}\right]_0^2=\frac{32}{5}$なので
\[
 T_4-A=\frac{53}{80},\qquad S_4-A=\frac1{60}.
\]
いずれも過大な近似である.\,$[0,2]$では$f''(x)=12x^2\leqq48$,\,$f^{(4)}(x)=24$だから,誤差上限は
\[
 |A-T_4|\leqq\frac{48\cdot2^3}{12\cdot4^2}=2,
 \qquad
 |A-S_4|\leqq\frac{24\cdot2^5}{180\cdot4^4}=\frac1{60}.
\]
台形公式では実際の誤差は上限より小さい.シンプソン公式では四階微分が一定なので,誤差を表す積分で絶対値を評価しても余分な差が生じず,上限と実際の誤差が一致する.

\paragraph{12.4}
(a) 分子,分母とも$0$へ近づき,\,$x\ne0$では分母の微分$2x$は$0$でない.左右それぞれでロピタルの定理を使う準備をする.微分後の商も$0$と$0$の形で,その分母の微分は$2\ne0$だから,もう一度適用できる.
\[
 \lim_{x\to0}\frac{\e^x-1}{2x}
 =\lim_{x\to0}\frac{\e^x}{2}=\frac12.
\]
この極限が存在するので,最初の商にも定理を適用でき,答えは$\frac12$である.

(b) $x>0$で両関数は微分可能であり,\,$\alpha>0$より
$\log x\to\infty$,\,$x^\alpha\to\infty$である.また分母の微分は$\alpha x^{\alpha-1}>0$なので
\[
 \frac{(\log x)'}{(x^\alpha)'}
 =\cfrac{\cfrac1x}{\alpha x^{\alpha-1}}
 =\frac1{\alpha x^\alpha}\longrightarrow0.
\]
よって求める極限は$0$である.

\paragraph{12.5}
$f(x)=\frac1{x^p(1-x)^q}$は$0<x<1$で正である.まず左端では
\[
 \frac{f(x)}{x^{-p}}=(1-x)^{-q}\longrightarrow1
 \qquad(x\to0+).
\]
極限比較により,左端で収束する条件は$\int_0^{\frac12}x^{-p}\dd x$と同じで,\,$p<1$である.
右端では
\[
 \frac{f(x)}{(1-x)^{-q}}=x^{-p}\longrightarrow1
 \qquad(x\to1-).
\]
$u=1-x$と置けば基準となる積分は$\int_0^{\frac12}u^{-q}\dd u$なので,右端の条件は$q<1$である.
両方の極限が独立に有限でなければならないから,必要十分条件は
\[
 \boxed{p<1\quad\text{かつ}\quad q<1}
\]
である.どちらかが満たされなければ,その端点の積分は$+\infty$へ発散する.

\paragraph{12.6}
まず有限の$R>\e$で$u=\log x$と置く.\,$\dd u=\frac1x\dd x$,下端は$u=1$,上端は$u=\log R$だから
\[
 \int_{\e}^R\frac1{x(\log x)^q}\dd x
 =\int_1^{\log R}u^{-q}\dd u.
\]
$R\to\infty$なら$\log R\to\infty$である.従って,これは無限大におけるべき型積分そのもので,収束条件は$q>1$である.その場合は
\[
 \int_{\e}^{\infty}\frac1{x(\log x)^q}\dd x
 =\lim_{R\to\infty}\frac{(\log R)^{1-q}-1}{1-q}
 =\boxed{\frac1{q-1}}.
\]
$q=1$では$\log\log R$,\,$q<1$では$(\log R)^{1-q}$が無限大に近づくため発散する.

\paragraph{12.7}
最初の積分では,左右を別に計算すると
\[
 \int_{-1}^{-\varepsilon}\frac1x\dd x=\log\varepsilon\to-\infty,
 \qquad
 \int_{\delta}^{1}\frac1x\dd x=-\log\delta\to+\infty.
\]
従って広義積分は発散する.一方,\,$\delta=\varepsilon$という対称な切り取り方では二つの和が常に$0$なので,対称的な主値は$0$である.

二番目の被積分関数を$g(x)=\frac{1-\cos x}{x}$とする.\,$1-\cos x=2\sin^2\frac x2$と$|\sin t|\leqq|t|$から
\[
 |g(x)|\leqq\frac{|x|}{2}\qquad(x\ne0)
\]
を得る.従って,左右それぞれで絶対値の積分が収束する.また$g(x)\to0$なので,\,$g(0)=0$と補えば連続関数になる.この$g$は奇関数だから
\[
 \int_{-1}^1\frac{1-\cos x}{x}\dd x=0.
\]
後者では各側が有限に収束するため,奇関数の対称性を使える.前者の$0$は,発散する二つの側を特定の切り取り方で相殺して得た主値である.

\paragraph{12.8}
被積分関数は偶関数である.\,$t=-u$による置換で
\[
 F(-x)=\int_0^{-x}\e^{-t^2}\dd t
 =-\int_0^x\e^{-u^2}\dd u=-F(x)
\]
となるので,\,$F$は奇関数である.基本定理より
\[
 F'(x)=\e^{-x^2}>0,\qquad F''(x)=-2x\e^{-x^2}.
\]
従って実数全体で狭義に単調増加し,\,$x<0$では下に凸,\,$x>0$では上に凸である.原点で凹凸が変わる.

$u\geqq0$に対して$\e^{-u}\leqq1$である.また$G(u)=\e^{-u}-1+u$と置くと,\,$G(0)=0$,\,$G'(u)=1-\e^{-u}\geqq0$なので$G(u)\geqq0$である.従って
\[
 1-u\leqq\e^{-u}\leqq1.
\]
$u=t^2$を代入し,\,$x\geqq0$として$0$から$x$まで積分すれば
\[
 \boxed{x-\frac{x^3}{3}\leqq F(x)\leqq x}
\]
を得る.

この上限$x$だけでは無限大での有界性は示せない.そこで$t\geqq1$では$t^2\geqq t$となることを使う.\,$x\geqq1$なら
\[
 F(x)=\int_0^1\e^{-t^2}\dd t+\int_1^x\e^{-t^2}\dd t
 \leqq1+\int_1^x\e^{-t}\dd t
 \leqq1+\e^{-1}.
\]
$F$は増加し,この一定値を超えないので,\,$x\to\infty$で有限の極限をもつ.
